% === Begin Label Data ===
% ID: 880
% 难度星级: 
% 标签: 
% 备注: 新定义
% 组卷引用次数: 0
% === End  Label Data ===

\begin{problem}{2020}{G}{北京卷}{21}{数列}
已知 $\{a_n\}$ 是无穷数列. 给出两个性质:

\circled{1} 对于 $\{a_n\}$ 中任意两项 $a_i, a_j (i > j)$, 在 $\{a_n\}$ 中都存在一项 $a_m$, 使得 $\displaystyle \frac{a_i^2}{a_j} = a_m$;

\circled{2} 对于 $\{a_n\}$ 中任意一项 $a_n (n \geqslant 3)$, 在 $\{a_n\}$ 中都存在两项 $a_k, a_l (k > l)$, 使得 $\displaystyle a_n = \frac{a_k^2}{a_l}$.

（I）若 $a_n = n (n = 1, 2, \dots)$, 判断数列 $\{a_n\}$ 是否满足性质 \circled{1}, 说明理由;

（II）若 $a_n = 2^{n-1} (n = 1, 2, \dots)$, 判断数列 $\{a_n\}$ 是否同时满足性质 \circled{1} 和性质 \circled{2}, 说明理由;

（III）若 $\{a_n\}$ 是递增数列, 且同时满足性质 \circled{1} 和性质 \circled{2}, 证明: $\{a_n\}$ 为等比数列.
\end{problem}

\begin{answer}
（I）不满足，理由：$\displaystyle \frac{a_3^2}{a_2} = \frac{9}{2} \notin \mathbb{N}^*$, 不存在一项 $a_m$ 使得 $\displaystyle \frac{a_3^2}{a_2} = a_m$.

（II）数列 $\{a_n\}$ 同时满足性质 \circled{1} 和性质 \circled{2}.

（III）证明见解析.
\end{answer}

\begin{solutions}
（I）不满足，理由：$\displaystyle \frac{a_3^2}{a_2} = \frac{9}{2} \notin \mathbb{N}^*$, 不存在一项 $a_m$ 使得 $\displaystyle \frac{a_3^2}{a_2} = a_m$.

（II）数列 $\{a_n\}$ 同时满足性质 \circled{1} 和性质 \circled{2},
理由：对于任意的 $i$ 和 $j$, 满足 $\displaystyle \frac{a_i^2}{a_j} = 2^{2i-j-1}$, 因为 $i \in \mathbb{N}^*, j \in \mathbb{N}^*$ 且 $i > j$, 所以 $2i - j \in \mathbb{N}^*$, 则必存在 $m = 2i - j$, 此时 $2^{m-1} \in \{a_i\}$ 且满足 $\displaystyle \frac{a_i^2}{a_j} = 2^{2i-j-1} = a_m$, 性质 \circled{1} 成立,
满足性质 \circled{2} 的方法一：对于任意的 $n$, 欲满足 $\displaystyle a_n = 2^{n-1} = \frac{a_k^2}{a_l} = 2^{2k-l-1}$, 满足 $n = 2k - l$ 即可,
因为 $k \in \mathbb{N}^*, l \in \mathbb{N}^*$, 且 $k > l$, 所以 $2k - l$ 可表示所有正整数,
所以必有一组 $k, l$ 使 $n = 2k - l$, 即满足 $\displaystyle a_n = \frac{a_k^2}{a_l}$, 性质 \circled{2} 成立.

满足性质 \circled{2} 的方法二：任取 $a_n = 2^{n-1} (n \geqslant 3)$, 总存在 $a_{n-1} = 2^{n-2}, a_{n-2} = 2^{n-3}$,
使得 $\displaystyle \frac{a_{n-1}^2}{a_{n-2}} = \frac{2^{2(n-2)}}{2^{n-3}} = 2^{n-1} = a_n$, 所以该数列具有性质 \circled{2}.

（III）假设数列 $\{a_n\}$ 各项均为正数, 观察目标, 记 $\displaystyle q = \frac{a_2}{a_1}$, 下面开始进行简单的数学实验:
$n = 3$ 时, 由性质 \circled{2}, 设 $\displaystyle a_3 = \frac{a_k^2}{a_l} = \frac{a_k}{a_l} \times a_k$, 因为 $3 > k > l \geqslant 1$, 所以 $k = 2, l = 1$.
即 $\displaystyle a_3 = \frac{a_2^2}{a_1} = a_1 q^2$, 所以 $a_1, a_2, a_3$ 成等比数列; (暂未用到性质 \circled{1})
$n = 4$ 时, 由性质 \circled{2}, 设 $\displaystyle a_4 = \frac{a_k^2}{a_l} = \frac{a_k}{a_l} \times a_k > a_k > a_l, 4 > k > l \geqslant 1$,
所以 $3 \geqslant k > l$, 进而 $\displaystyle a_4 = \frac{a_k^2}{a_l} = a_1 q^{2k-l-1}$, 记 $n_4 = 2k - l - 1 \in \mathbb{N}^*$.
由性质 \circled{1}, 存在正整数 $m$, 使得 $\displaystyle a_m = \frac{a_4^2}{a_2} = q \times a_3 = a_1 q^3$, 由 $q > 1$, 得 $a_m > a_3$.
若 $a_4 = a_1 q^{n_4} < a_1 q^3$, 再由 $a_3 = a_1 q^2 < a_4$, 得 $2 < n_4 < 3$, 与 $n_4 \in \mathbb{N}^*$ 矛盾;
若 $a_4 > a_1 q^3 = a_m$, 则 $a_3 < a_m < a_4$, 与数列 $\{a_n\}$ 递增矛盾,
所以 $a_4 = a_1 q^3$.

下用数学归纳法完成证明:
假设当 $1 \leqslant n \leqslant r$ 时, $a_n = a_1 q^{n-1}$ 都成立.
探究 $a_{r+1} = a_1 q^r$
由性质 \circled{2}, 存在正整数 $k, l (k > l)$, 使得 $\displaystyle a_{r+1} = \frac{a_k}{a_l} \times a_k > a_k > a_l$, 得 $r + 1 > k > l$,
所以 $r \geqslant k > l$, 进而 $\displaystyle a_{r+1} = \frac{a_k^2}{a_l} = a_1 q^{2k-l-1}$, 记 $n_{r+1} = 2k - l - 1 \in \mathbb{N}^*$.
由性质 \circled{1}, 存在正整数 $h$, 使得 $\displaystyle a_h = \frac{a_r^2}{a_{r-1}} = q a_r = a_1 q^r$, 由 $q > 1$, 得 $a_h > a_r$.
下面我们证明 $a_{r+1} = a_1 q^r$,
若 $a_{r+1} = a_1 q^{n_r} < a_1 q^r$, 再由 $a_r = a_1 q^{r-1} < a_{r+1}$, 得 $r < n_{r+1} < r + 1$, 与 $n_{r+1} \in \mathbb{N}^*$ 矛盾,
若 $a_{r+1} > a_1 q^r = a_h$, 则 $a_r < a_h < a_{r+1}$, 与数列 $\{a_n\}$ 递增矛盾, 所以 $a_{r+1} = a_1 q^r$.
由数学归纳法知, 对任意 $n \in \mathbb{N}^*$, 都有 $a_n = a_1 q_{n-1}$, 所以数列 $\{a_n\}$ 是等比数列.
若数列全为负, 则同理可得;

再考虑数列不可能同时存在正负项:
证明方法一 反证法：假设这个递增数列先负后正,
那么必有一项 $a_l$ 绝对值最小或者有 $a_l$ 与 $a_{l+1}$ 同时取得绝对值最小,
如仅有一项 $a_l$ 绝对值最小, 此时必有一项 $\displaystyle a_m = \frac{a_l^2}{a_j}$, 此时 $|a_m| < |a_l|$ 矛盾;
如有两项 $a_l$ 与 $a_{l+1}$ 同时取得绝对值最小值, 那么必有 $\displaystyle a_m = \frac{a_l^2}{a_{l+1}}$, 此时 $|a_m| = |a_l|$, 矛盾;
证毕.

证明方法二：由性质 \circled{1}, $i, j \in \mathbb{N}^* (i > j)$, $\displaystyle \frac{a_i^2}{a_j}$ 都是数列 $\{a_n\}$ 中的项, 所以 $a_j \neq 0 (j \in \mathbb{N}^*)$.
假设 $a_p < 0 < a_{p+1}$, 因为若 $\{a_n\}$ 是递增数列, 所以数列 $\{a_n\}$ 中共有 $p$ 个负项.
另一方面, 由性质 \circled{1}, $\displaystyle \frac{a_{p+h}^2}{a_p} (h \in \mathbb{N}^*)$ 都是数列 $\{a_n\}$ 中的项,
且 $\displaystyle 0 > \frac{a_{p+1}^2}{a_p} > \frac{a_{p+2}^2}{a_p} > \dots > \frac{a_{p+h}^2}{a_p} > \dots$, 与数列 $\{a_n\}$ 中共有 $p$ 个负项矛盾.
所以同时满足的数列恒正或者恒负.
\end{solutions}